\section{The microsimulation stage}\label{sec:roadmap}

The public-data envelopes answer ``roughly how much''; they cannot answer
``who receives it, net of what, at what net Exchequer cost.'' Those are
tax--benefit questions, and the second stage answers them with the same
architecture as the companion study: FRS 2024--25 households in
PolicyEngine UK, with each arm implemented as a hypothetical statutory
rule.

Four interactions make the net numbers differ from the gross ones, and
each is mechanical once the arms are in the rules engine. First, wage
insurance paid as taxable earnings is taxed at the recipient's marginal
rate and tapered against Universal Credit at 55 pence in the pound: a
household on UC keeps well under half of a gross wage-insurance pound, so
the net Exchequer cost of Arm A is far below its gross cost---and its
distributional incidence tilts away from the means-tested poor, an effect
with no US analogue and, to my knowledge, no existing quantification
anywhere. Second, the UI arm substitutes for UC pound-for-pound where
recipients are entitled to both, so its incremental cost depends on the
joint distribution of contribution histories and household means. Third,
the redundancy-disregard arm's cost is precisely the UC entitlement the
frozen capital thresholds currently extinguish---computable only from the
joint capital--redundancy distribution, for which the microsimulation
stage will combine FRS capital data with statutory redundancy-pay
formulas applied to observed tenure and age. Fourth, every arm's poverty
and decile incidence runs through household composition, which the flow
arithmetic of Section~\ref{sec:method} cannot see.

The displacement population for this stage is the broad-eligibility one:
all FRS manufacturing employees (roughly 1,100 records) carrying
LFS-estimated redundancy risk, evaluated by probability-weighted
enumeration---each worker's displacement counterfactual simulated once
and weighted by their risk---rather than by Monte Carlo assignment. This
matters beyond convenience: the companion study's trade-calibrated
scenario concentrates its loss on so few survey records that assignment
noise dominates subgroup statistics, whereas a policy costing over the
full at-risk population uses two orders of magnitude more survey support
and needs no assignment draws at all. The trade-calibrated scenario
returns as the narrow-eligibility bounding case, not the workhorse. The
first implementation step---simulating New Style JSA (and its proposed UI
successor) in the displacement margin---is shared with the companion
study's own revision programme, where the unmodelled JSA channel is the
referee's leading demand; one build serves both papers.

Behaviour enters the second stage the same way it enters the first:
as bracketed scenarios, now with more structure. The
\citet{hyman2024} employment effects imply shorter non-employment spells
under Arm A; implemented as duration scenarios over the enumeration
population, they bound how much of the gross cost the Exchequer recoups
through earlier re-employment---the UK answer to whether wage insurance
can self-finance once the taper clawback, which the US setting lacks, is
taken into account.
